\subsection{The exact mixed-chart image of the geometric measure}
\label{subsec:cm-measure-image}

Write \(w_K=|O_K^\times|\), put \(U=(\mathbb Z_p^\times)^2\subset V\), and let
\(j:U\hookrightarrow V\) be its inclusion.
The ray-class identification is
\(G=U/O_K^\times\).  For the auxiliary ideal
\(\mathfrak a=(c)\), let \(\Phi_c\) denote the map of Lemma~\ref{lem:cm-mixed-measure-embedding} composed with \(A\to\overline A\), and let \(\mu_c\) be its normalized
\(O_K^\times\)-invariant lift of the mixed-unit
measure \(\Phi_c\) to \(U\): its mass on each finite
unit orbit is the quotient mass divided by \(w_K\).
This is integral since \(p\nmid w_K\).
Let \(P,Q\) act on \(V\) by multiplication by
\(\pi,\bar\pi\), and put
\[
 \mathcal L=(1-P_*)(1-Q_*),\qquad
 D_c=c^2-[c^{-1}]_*.
\]
Here all starred maps denote pushforward of bounded
measures, and \(PQ=p\,\operatorname{id}_V\).

\begin{lemma}\label{lem:cm-measure-euler-image}
The actual geometric measure satisfies
\[
 w_Kj_*\mu_c=D_c\mathcal L\alpha,\qquad
 w_Kj_*\mu_c=D_cj_*(\alpha|_U),\qquad
 \mathcal L\alpha=j_*(\alpha|_U).
\]
Its restriction \(\alpha|_U\) has kernel
\(N_{\mathrm{ell}}\), and uniquely determines
\(\alpha\) with its zero-total-mass normalization.
Under finite \(\mathcal X_2^+\), this restriction
embeds \(A\) as a full-rank closed submodule of
\(\mathbb Z_p[[U]]^{O_K^\times}(1)\).
After identifying this target with
\(\mathbb Z_p[[G]](1)\) and absorbing the twist,
the image is an ideal containing a nonzerodivisor.
\end{lemma}

\begin{proof}
The \((r,s)\)-moment of \(\mu_c\) is
\(\kappa_{(r+1,s+1),\mathfrak a}\).  Rearranging
\cite[Proposition 4.3, pp.~17--18]{IshiiSoule}
and using \cite[Theorem 2.6, p.~8]{Ishii}
gives exactly the first identity on every admissible
positive-degree monomial.  Inadmissible moments
vanish by \(O_K^\times\)-invariance.  The constant
moment of \(\mu_c\) is zero: it is the Kummer value
of \(N_{L_n/K}(u_n)\), a global unit of \(K\), hence
a root of unity with all \(p\)-power roots in
\(L_\infty\).  The constant moment of \(\alpha\)
is zero by its source formula, and the odd moments
vanish by the action of \(-1\).
Polynomial moments determine bounded measures,
so the first identity holds integrally.

Neither \(P(V)\) nor \(Q(V)\) meets \(U\).
Restriction to \(U\), which commutes with \(D_c\),
therefore gives the second identity.
The \((r,s)\)-moment multiplier of \(D_c\) is
\(c^2-c^{-r-s}\ne0\), including degree zero.
Thus \(D_c\) is injective, and cancellation gives
the third identity.

The kernel of \(\mathcal L\) on all bounded
measures is precisely \(\mathbb Z_p\delta_0\).
Indeed its positive-degree moment multipliers
are
\[
 (1-i^{(r,s)}(\pi))(1-i^{(r,s)}(\bar\pi)),
\]
and both algebraic numbers occurring here have
complex absolute value \(p^{(r+s)/2}>1\).
A measure killed by \(\mathcal L\) consequently
has all positive-degree moments zero and is its
total mass times \(\delta_0\).
This proves uniqueness with zero total mass.
The second identity and injectivity of \(D_c\)
identify \(\ker(\alpha|_U)\) with
\(\ker\Phi_c=N_{\mathrm{ell}}\).

Under finite \(\mathcal X_2^+\), the latter kernel
vanishes.  The exact rank formula
\cite[Theorem 4.4, author version p.~27]{VenjakobIwasawa},
together with weak Leopoldt, gives
\[
 \operatorname{rank}_{\Lambda_2}A
     =r_2(L_1)=[L_1:K]=\frac{(p-1)^2}{w_K}.
\]
The maximal pro-\(p\) group comparison used here
is the \((p,S)\)-closed comparison at the start of
\cite[Section 4, author version p.~26]{VenjakobIwasawa}.
The unit-chart target has this same rank.
Equivariance makes the image an ideal after the
cyclotomic twist is absorbed.  The tame
decomposition of \(\mathbb Z_p[[G]]\) is a product
of copies of \(\Lambda_2\), and full rank makes
each ideal component nonzero.  A choice of a
nonzero element in every component gives the
claimed nonzerodivisor.
\end{proof}

These identities retain the radial and axis constraints
on the actual image.  In particular, for \(i,j\geq0\),
\[
 \alpha|_{P^iQ^jU}=(P^iQ^j)_*(\alpha|_U).
\]
This follows by restricting the Euler identity
successively to the disjoint interior shells.
No infinite geometric-series expansion at the axes
is required.  In the class-two question, the remaining
map is still
\[
 \mathcal J\circ(\alpha\widehat\otimes\alpha)
              \bigm|_{\widehat{\bigwedge}^{\,2}A}.
\]
The separate-factor operators \(\mathcal L\otimes1\)
and \(1\otimes\mathcal L\) do not commute with its
addition shears.  The displayed ideal description
therefore retains, rather than removes, this exact
tensor-image constraint.
